\documentclass[12pt,a4paper]{report}

\usepackage[utf8]{inputenc}
\usepackage[T5]{fontenc}
\usepackage[vietnamese]{babel}
\usepackage{graphicx}
\usepackage{hyperref}
\usepackage{geometry}
\usepackage{titlesec}
\usepackage{tocloft}

\geometry{a4paper, left=3cm, right=2cm, top=2cm, bottom=2cm}

\titleformat{\chapter}[display]
  {\normalfont\huge\bfseries\centering}{\chaptertitlename\ \thechapter}{20pt}{\Huge}
\titlespacing*{\chapter}{0pt}{-50pt}{40pt}

\hypersetup{
    colorlinks=true,
    linkcolor=black,
    filecolor=black,      
    urlcolor=blue,
    pdftitle={Báo cáo Đồ án Big Data - SOC},
}

\begin{document}

\begin{titlepage}
    \begin{center}
        \textbf{\Large ĐẠI HỌC SÀI GÒN} \\[0.5cm]
        \textbf{\large KHOA TOÁN - ỨNG DỤNG} \\[3cm]

        \textbf{\Large BÁO CÁO ĐỒ ÁN CUỐI KÌ KẾT THÚC HỌC PHẦN} \\[0.5cm]
        \textbf{\large Bộ Môn: Xử lý dữ liệu lớn (Big Data)} \\[2cm]

        {\LARGE \textbf{CHỦ ĐỀ: XÂY DỰNG HỆ THỐNG SOC PHÂN TÍCH DỮ LIỆU LỚN VỚI KAFKA VÀ SPARK}}\\[3cm]

        \begin{flushleft}
            \large
            \textbf{Giảng viên:} Vũ Ngọc Thanh Sang \\[0.5cm]
            \textbf{Sinh Viên:} Nguyễn Trương Cao Sơn \\[0.5cm]
            \textbf{Lớp:} DDU1231 
        \end{flushleft}

        \vfill
        {\large TP Hồ Chí Minh, năm 2024}
    \end{center}
\end{titlepage}

\pagenumbering{roman}

\tableofcontents
\newpage

\listoffigures
\newpage

\pagenumbering{arabic}

\chapter*{LỜI MỞ ĐẦU}
\addcontentsline{toc}{chapter}{LỜI MỞ ĐẦU}
(Viết lời mở đầu ở đây...)

\newpage

\chapter{TỔNG QUAN VỀ ĐỀ TÀI}

\section{Giới thiệu về đề tài}
\subsection{Bối cảnh nghiên cứu}
(Nội dung...)

\subsection{Lý do chọn đề tài}
(Nội dung...)

\section{Mục tiêu và phạm vi nghiên cứu}
\subsection{Mục tiêu nghiên cứu}
(Nội dung...)

\subsection{Phạm vi giới hạn của đồ án}
(Nội dung...)

\section{Ý nghĩa thực tiễn}
(Nội dung...)

\chapter{CƠ SỞ LÝ THUYẾT VÀ CÔNG NGHỆ}

\section{Khái quát về Xử lý dữ liệu lớn (Big Data)}
\subsection{Đặc trưng của dữ liệu lớn (Mô hình 3V/5V)}
(Nội dung...)

\subsection{Xử lý theo lô (Batch) và Xử lý luồng (Stream Processing)}
(Nội dung...)

\section{Nền tảng phân phối thông điệp Apache Kafka}
\subsection{Kiến trúc tổng quan của Kafka}
(Nội dung...)

\subsection{Cơ chế lưu trữ và hoạt động (Topic, Partition)}
(Nội dung...)

\section{Khung tính toán phân tán Apache Spark}
\subsection{Giới thiệu về Apache Spark và xử lý In-Memory}
(Nội dung...)

\subsection{Mô hình Spark Structured Streaming}
(Nội dung...)

\section{Ứng dụng trong An toàn thông tin (Hệ thống SIEM)}
\subsection{Khái niệm về hệ thống SIEM}
(Nội dung...)

\subsection{Kỹ thuật tấn công Brute Force và cơ chế phát hiện}
(Nội dung...)

\chapter{ỨNG DỤNG VÀ TRIỂN KHAI BÀI TOÁN THỰC TẾ}

\section{Phương pháp nghiên cứu và kiến trúc hệ thống}
\subsection{Sơ đồ khối kiến trúc tổng thể}
(Nội dung...)

\subsection{Lựa chọn công nghệ triển khai (Docker Containerization)}
(Nội dung...)

\section{Thu thập và giả lập dữ liệu (Data Ingestion)}
\subsection{Xây dựng kịch bản mạng Botnet tấn công}
(Nội dung...)

\subsection{Cấu hình Kafka Producer}
(Nội dung...)

\section{Xử lý và phân tích dữ liệu (Data Processing)}
\subsection{Xử lý luồng dữ liệu với Spark Structured Streaming}
(Nội dung...)

\subsection{Thuật toán phát hiện bất thường và trích xuất Blacklist}
(Nội dung...)

\chapter{KẾT QUẢ THỰC NGHIỆM VÀ ĐÁNH GIÁ}

\section{Kết quả chạy thực nghiệm}
\subsection{Giao diện sinh dữ liệu log}
(Nội dung...)

\subsection{Bảng cảnh báo thời gian thực và File kết quả}
(Nội dung...)

\section{Đưa ra nhận xét và đánh giá}
\subsection{Ưu điểm của hệ thống}
(Nội dung...)

\subsection{Hạn chế và hướng phát triển}
(Nội dung...)

\chapter*{KẾT LUẬN}
\addcontentsline{toc}{chapter}{KẾT LUẬN}
(Nội dung kết luận...)

\chapter*{LỜI CẢM ƠN}
\addcontentsline{toc}{chapter}{LỜI CẢM ƠN}
(Nội dung lời cảm ơn...)

\chapter*{TÀI LIỆU THAM KHẢO}
\addcontentsline{toc}{chapter}{TÀI LIỆU THAM KHẢO}
(Liệt kê tài liệu...)

\end{document}